% Source: https://artofproblemsolving.com/community/c3h1425731_absolute_value_equation
% Stable identity: aops-topic-1425731-q01
The expressions inside the absolute values change sign at $x=-1$, $x=-\dfrac{1}{3}$ and $x=5$.
\begin{itemize}
\item $x\ge5$: $(3x+1)+(4x-20)+(x+1)=8x-18=40$, so $x=\dfrac{29}{4}$, which is $\ge5$. Valid.
\item $-\dfrac{1}{3}\le x<5$: $(3x+1)+(20-4x)+(x+1)=22\ne40$. No solution.
\item $-1\le x<-\dfrac{1}{3}$: $(-3x-1)+(20-4x)+(x+1)=20-6x=40$, so $x=-\dfrac{10}{3}$, not in this interval.
\item $x<-1$: $(-3x-1)+(20-4x)+(-x-1)=18-8x=40$, so $x=-\dfrac{11}{4}$, which is $<-1$. Valid.
\end{itemize}
The solutions are $\boxed{x=-\dfrac{11}{4}\text{ and }x=\dfrac{29}{4}}$.
